% Isolated post-CP36 candidate. No canonical or checkpoint input is changed.
\subsection{Actual evaluation fibres for a polarization of type \((1,5)\)}

Throughout this subsection, \(k\) is algebraically closed of
characteristic two, \(A/k\) is an abelian surface, and \(L\) is
an ample globally generated line bundle of type \((1,5)\).
Thus \(W=H^0(A,L)\) has dimension five and \(L^2=10\).
Write
\[
 F_A:A\longrightarrow B=A^{(2)},\qquad
 \mathcal G=F_{A*}L,\qquad H=L^{(2)}.
 \tag{EF5.1}
\]
Here \(F_A\) is relative Frobenius, whereas \(H\) is a line
bundle on \(B\); the canonical identity is
\(F_A^*H\simeq L^{\otimes2}\).

\begin{lemma}[The actual kernel map and its Cartier image]
\label{lem:DE-five-evaluation-kernel}
\status{Candidate; full independent proof audits pending.}
There is a nonempty open \(U_B\subset B\) on which evaluation
\[
              W\otimes\mathcal O_B\longrightarrow\mathcal G
                                                        \tag{EF5.2}
\]
is surjective. Its kernel on \(U_B\) is a line subbundle and
defines a morphism \(\psi_B:U_B\to\mathbf P(W)\).
For \(U_A=F_A^{-1}(U_B)\), put
\(\psi=\psi_B\circ F_A:U_A\to\mathbf P(W)\).
At a geometric point \(z\in U_A\), this is exactly the line
\[
 \psi(z)=\ker\bigl(W\longrightarrow
          H^0(F_A^{-1}(F_A(z)),L|_{F_A^{-1}(F_A(z))})\bigr).
 \tag{EF5.3}
\]

If \(A\) is nonordinary, fix a compatible level-five theta
basis and let \(B_y\), \(\sigma(y)\) and
\(Q(y)=\sum_i y_i\sigma_i(y)^2\) be its actual dual Cartier
matrix, maximal-Pfaffian vector and polar polynomial. For the
complete coefficient morphism
\(\varphi_H:B\to\mathbf P(H^0(B,H)^*)\), one has on \(U_B\)
\[
 \operatorname{rank}B_{\varphi_H(x)}=4,
 \qquad \psi_B(x)=[\sigma(\varphi_H(x))].
 \tag{EF5.4}
\]
In particular the entire actual coefficient image satisfies
\[
                     Q\circ\varphi_H=0.
 \tag{EF5.5}
\]
No embedding property of \(\varphi_H\) is asserted.
\end{lemma}
\begin{proof}
Proposition~\ref{prop:DE-prime-to-six-evaluation-four} applies
with \(N=5\), since \(\gcd(5,6)=1\). It proves generic
surjectivity of (EF5.2), including when \(A\) is ordinary.
The bundle \(\mathcal G\) has rank four, so its kernel on the
surjective open is a line bundle. Finite flat base change for
relative Frobenius identifies its fibre with (EF5.3).
On the completed local ring at \(z\), parameters \(u,v\)
identify the Frobenius fibre ideal with \((u^2,v^2)\).

In the nonordinary case the relative Cartier pairing
\[
                  \mathcal G\otimes\mathcal G\longrightarrow H
 \tag{EF5.6}
\]
is perfect and alternating. The value \(B_{\varphi_H(x)}\)
is precisely the pullback of this pairing on the four-dimensional
fibre \(\mathcal G_x\) along the surjective evaluation map.
Consequently its radical is the kernel of evaluation and its
rank is four. The Pfaffian radical formula gives (EF5.4).

For completeness, (EF5.5) is an actual evaluation identity.
Let \(x=F_A(z)\), choose compatible local trivializations of
\(L\) and \(H\), and write the basis of \(W\) as
\(e_0,\ldots,e_4\). The evaluation vector is
\(y_i=e_i^{(2)}(x)=e_i(z)^2\). Since the section
\(s=\sum_i\sigma_i(y)e_i\) belongs to the evaluation kernel,
it vanishes on the whole Frobenius fibre and hence at \(z\).
Therefore
\[
 Q(y)=\sum_i e_i(z)^2\sigma_i(y)^2
       =\left(\sum_i\sigma_i(y)e_i(z)\right)^2=0.
 \tag{EF5.7}
\]
This is independent of the chosen trivializations. It holds on
the dense open \(U_B\); since \(B\) is integral, it is the
zero section of \(H^{\otimes5}\) everywhere. In this argument
the kernel section is allowed to depend on \(x\); no fixed
global kernel section is inferred.
\end{proof}

\begin{proposition}[The elliptic-degree-two dichotomy]
\label{prop:DE-five-evaluation-fibre-dichotomy}
\status{Candidate; full independent proof audits pending.}
The map \(\psi\) has positive-dimensional generic fibres if
and only if \(A\) contains an elliptic curve \(E\) with
\(L\cdot E=2\). Equivalently, \(\psi\) is generically finite
onto its image if and only if there is no such elliptic curve.

More precisely, the closure \(D\) of every curve component of
a fibre of \(\psi\) is a smooth elliptic translate and satisfies
\(L\cdot D=2\). If \([s]\) is that fibre's section, then
\[
               \operatorname{div}(s)=2D+R,
        \qquad R^2=2,\quad L\cdot R=6,
 \tag{EF5.8}
\]
and \(\mathcal O_A(R)\) is an ample principal polarization.

Conversely, given an elliptic subgroup \(E\subset A\) of
\(L\)-degree two, let \(\pi:A\to C=A/E\) be the quotient
and \(E_x=\pi^{-1}(x)\). The unique section, up to scalar, of
the ample principal line \(L(-2E_x)\), multiplied by the
canonical section of \(\mathcal O_A(2E_x)\), defines a
morphism
\[
          \gamma:C\longrightarrow\mathbf P(W),
        \qquad \psi=\gamma\circ\pi\quad\hbox{on }U_A.
 \tag{EF5.9}
\]
The map \(\gamma\) is nonconstant and injective on geometric
points. No separability of \(\gamma\) is asserted.
\end{proposition}
\begin{proof}
First let \(D\) be the closure of a curve in a fibre of
\(\psi\), with constant kernel section \([s]\).
Equation (EF5.3) says that \(s\) lies in the Frobenius-fibre
ideal at every point in a dense open of \(D\). This implies
\[
                         \operatorname{div}(s)\ge2D.
 \tag{EF5.10}
\]
Indeed \(D\) is generically smooth because \(k\) is perfect.
If its multiplicity in \(\operatorname{div}(s)\) were one,
then at a general smooth point of \(D\) a local equation for
\(s\) would be a local parameter times a unit. Its linear
term would be nonzero, whereas the ideal \((u^2,v^2)\) is
contained in the square of the maximal ideal. This contradicts
(EF5.3). This argument concerns the actual fibre ideal, not only
the vanishing of a differential at a special point.

Write \(\operatorname{div}(s)=2D+R\) with \(R\) effective.
The residual divisor is nonzero: otherwise \(L^2=4D^2\),
which is impossible for \(L^2=10\). Every effective divisor
on an abelian surface is nef. To see this, translate it so that
it contains no chosen curve, then take the resulting nonnegative
proper intersection number. In particular \(D^2,R^2\ge0\),
and self-intersections are even by Riemann--Roch.

If \(D^2>0\), then \(D^2\ge2\), and \(D\) is ample.
For the latter assertion, Hodge index and nefness show that a
nonzero effective curve cannot have zero intersection with a
nef divisor of positive square; Nakai's criterion then applies.
Since \(R\ne0\), this gives \(D\cdot R\ge1\), and hence
\[
       10=L^2=4D^2+4D\cdot R+R^2\ge8+4=12,
\]
a contradiction. Thus \(D^2=0\).

Adjunction gives \(p_a(D)=1\). An abelian variety contains no
rational curve, so the normalization has genus at least one.
It follows that \(D\) is smooth of genus one. After translating
one of its points to zero, its inclusion in \(A\) is a group
homomorphism; thus \(D\) is an elliptic translate. Global
generation restricts to \(L|_D\), whose positive degree cannot
be one on an elliptic curve. Consequently \(L\cdot D\ge2\).
As \(D^2=0\), the intersection identity becomes
\[
                    10=4L\cdot D+R^2.
 \tag{EF5.11}
\]
It forces \(L\cdot D=2\) and \(R^2=2\). Also
\(L\cdot R=10-2L\cdot D=6\). As above, an effective nef
divisor with positive square is ample; its Euler characteristic
is one, so it is principal. This proves (EF5.8).

The image of \(\psi\) cannot be a point: a constant nonzero
section would vanish on the Frobenius fibres of a dense open
of \(A\), hence would vanish identically. If \(\psi\) is
not generically finite, its image therefore has dimension one;
a general fibre supplies a curve \(D\) just considered.
This proves one direction of the dichotomy.

For the converse, every fibre \(E_x\) has the same numerical
class as \(E\). Put \(R_x=L(-2E_x)\). Then
\[
                     R_x^2=2,\qquad R_x\cdot L=6.
 \tag{EF5.12}
\]
These inequalities make \(R_x\) ample. Here is a numerical
justification that does not assume effectiveness in advance.
In the real Neron--Severi space, Hodge index makes the intersection
form have signature \((1,\rho-1)\). The two inequalities
(EF5.12) put \(R_x\) in the component of the positive cone
containing \(L\). Every nonzero effective curve is nef, has
nonnegative square and has positive intersection with \(L\);
it is in the closure of that component. Thus its intersection
with \(R_x\) is strictly positive. Nakai's criterion proves
the assertion. Abelian ample-line vanishing and Riemann--Roch
now give \(h^0(R_x)=\chi(R_x)=1\).

The sections in (EF5.9) form an actual algebraic family. On
\(A\times C\), let \(\mathcal D\) be the Cartier divisor
\(\{(z,x):\pi(z)=x\}\), and put
\[
 \mathcal R=\operatorname{pr}_A^*L(-2\mathcal D).
\]
Its restrictions to the fibres over \(C\) are the ample lines
\(R_x\), with \(h^0=1\) and all higher cohomology zero.
Cohomology and base change therefore make
\(\operatorname{pr}_{C*}\mathcal R\) a line bundle.
Its evaluation map, multiplied by the canonical section of
\(2\mathcal D\), gives a line subbundle of
\(W\otimes\mathcal O_C\), and hence the morphism \(\gamma\).

At a point \(z\in E_x\), the resulting section is divisible
by the square of a local equation for \(E_x\). It therefore
vanishes on the Frobenius fibre at \(z\). On \(U_A\) the
evaluation kernel is one-dimensional, which proves (EF5.9).
If \(x\ne x'\) and \(\gamma(x)=\gamma(x')\), that common
section would have divisor at least \(2E_x+2E_{x'}\). The
remaining effective divisor would have numerical class
\(L-4E\), whose square is \(10-8\cdot2=-6\), impossible
for an effective divisor on an abelian surface. Thus \(\gamma\)
is injective on geometric points and nonconstant. Its image
has dimension one, completing the converse.
\end{proof}

\begin{remark}[What the fibre dichotomy does not decide]
\label{rem:DE-five-evaluation-fibre-scope}
\status{Candidate scope; no exceptional-coefficient exclusion claimed.}
The proposition is an actual geometric classification for a
globally generated polarization of type \((1,5)\), and its
dichotomy does not require \(A\) to be nonordinary. Nonordinarity
is used for the alternating Cartier and polar assertions in
Lemma~\ref{lem:DE-five-evaluation-kernel}. Neither generic
generation of \(F_{A*}L\) nor generic finiteness of \(\psi\)
is upgraded to pointwise generation or an embedding.

In particular, this does not exclude or realize the exceptional
prime-five Heisenberg ratios. A divisor contracted by the ambient
birational Pfaffian map can still map a surface inside it
generically finitely onto an inverse-base surface. Thus even
the generically finite side of the dichotomy is not an
exceptional-coefficient exclusion. No common-maximal-rank wall,
inertia, generic Lang connectivity or full monodromy assertion
follows here. The proof uses the already audited odd-evaluation,
theta and abelian-divisor arguments and elementary intersection
theory; no new numerical computation is claimed.
\end{remark}
